\chapter{No Escape to Infinity}\label{ch:no-escape-to-infinity}

A key difficulty in min--max theory on \emph{complete} manifolds (with or without boundary) is ensuring that the min--max sequence does not ``escape to infinity''---that is, that slices with large area do not concentrate all their mass arbitrarily far from the good set $U$. As observed in \cite{CL20}, this is fundamentally a problem caused by the \emph{completeness} of the manifold: in a compact manifold, such escape is topologically impossible, but in a complete non-compact setting, one must prove that the min--max construction forces slices to remain localized near $U$.

This chapter establishes a quantitative non-escape result: any good sweepout must have at least one slice carrying a definite amount of area inside $U$. The free-boundary setting introduces additional technical complications in the proof (particularly in handling boundary collars and relative isoperimetric estimates), but the fundamental difficulty remains the completeness of the ambient manifold.

The proof proceeds by contradiction. We assume such a sweepout exists and use the nested sweepout structure from Chapter~\ref{ch:nested-sweepout} to construct a competing sweepout with strictly smaller width, contradicting the definition of $W_g(U)$.

\section{Statement of the Main Result}\label{sec:non-escape-statement}

The main result of this chapter is the following proposition, which provides a quantitative lower
bound on the area that must remain inside $U$.

\begin{proposition}\label{prop:non-escape}
    For every good set $U$, there exists a positive constant $\eps(U)$ which depends only on $U$ such that the following holds: For every good sweepout $\br{\Gamma_t}$ of $U$ with associated family of open sets $\br{\Omega_t}$, there exists $t' \in [0, 1]$ such that
    \begin{equation*}
        \mathcal{H}^n(\Gamma_{t'}) \geq W_g(U),
    \end{equation*}
    and such that $\mathcal{H}^n(\Gamma_{t'}\cap cl(U)) \geq \eps(U)$.
\end{proposition}

The proof is by contradiction. We assume that Proposition~\ref{prop:non-escape} does not hold, i.e. we can construct a good sweepout with volume of hypersurfaces strictly less than $W_g(U)$. The main tool in the proof is the structure of the nested sweepout from Proposition~\ref{prop:nested-sweepout}.

\section{Pushing Sets Outside Good Sets}\label{sec:pushing-outside}

The first technical ingredient is a lemma that allows us to deform a set with small intersection with $U$ so that it avoids $U$ entirely, while controlling the area increase.

Let $U$ be a good set.
\begin{lemma}\label{lem:nested-sweepout-small-intersection}
    There exist $\eps(U) > 0$, $\eps_0(U) > 0$ and $\eps_1(U) > 0$ such that for any open sets $\Omega'$ the following holds:
    \begin{enumerate}
        \item $\max\br{\eps, \eps_1} < \mathbf{M}(\partial_IU)/10$;
        \item If $\eps_0 < \mathcal{H}^{n+1}(\Omega'\cap U) < \mathcal{H}^{n+1}(U) - \eps_0$, then $\mathbf{M}(\partial_I\Omega'\cap U) > 2\eps$;
        \item 
        \begin{enumerate}
            \item[(A)] If $\mathcal{H}^{n+1}(\Omega'\cap U) < 2\eps_0$, then there exists a family of open sets $\br{\Xi_t}$ with 
            \begin{equation*}
                \begin{split}
                    &\Xi_0 = \Omega';\\
                    &\Xi_t\backslash N_{\eps_1}(U) = \Omega'\backslash N_{\eps_1}(U);\\
                    &\Xi_1\cap U = \varnothing\\
                    &\mathbf{M}(\partial_I\Xi_t) < \mathbf{M}(\partial_I\Omega') + \mathbf{M}(\partial_IU) + \eps_1.
                \end{split}
            \end{equation*}
            \item[(B)] If $\mathcal{H}^{n+1}(\Omega'\cap U) > \mathcal{H}^{n+1}(U) - 2\eps_0$, then there exists a family of open sets $\br{\Xi_t}$ with
            \begin{equation*}
                \begin{split}
                    &\Xi_0 = \Omega';\\
                    &\Xi_t\backslash N_{\eps_1}(U) = \Omega'\backslash N_{\eps_1}(U);\\
                    &\Xi_1\cap U = U\\
                    &\mathbf{M}(\partial_I\Xi_t) < \mathbf{M}(\partial_I\Omega') + \mathbf{M}(\partial_IU) + \eps_1.
                \end{split}
            \end{equation*}
        \end{enumerate}
    \end{enumerate}
    \begin{proof}
        Pick some $\eps_1 \in (0, \mathcal{H}^{n}(\partial_IU)/10)$. 
        
        We will show that for all sufficiently small $\eps_0$ (with the choice of $\eps_0$ depending on $\eps_1$) statement (3) holds;  

        We will show that for all sufficiently small $\eps$ (with the choice of $\eps$ depending on $\eps_0$) statement (2) holds; the argument relies on the convexity properties of the isoperimetric profile of $\overline{U}$.

        We will prove Statement(3)(A). And the statement (3)(B) would be similar.

        Let $r_0 > 0$ be sufficiently small, so that every ball $B$ of radius $r \in (0, r_0]$ centered at a point in $cl(U)$ is $2$-biLipschitz diffeomorphic to a ball of the same radius in the Euclidean space.

        Choose a covering $\br{B_i}$ of $cl(U)$ by $N$ balls of radius $r_0$, so that concentric balls of radius $\frac{r_0}{2}$, denoted by $\frac{1}{2}B_i$, still cover $cl(U)$. Set
        \begin{equation}\label{eq:value-of-eps0}
            \eps_0 := \min\br{\frac{\eps_1 r_0}{20 N}, \brac{\frac{\eps_1}{2^{n + 2}C(n)N}}^{\frac{n + 1}{n}}},
        \end{equation}
        where $C(n)$ is the constant from Lemma~\ref{lem:morse-separate}. Using the relative coarea formula (Corollary~\ref{cor:coarea-filling}), we may choose a covering $\br{B_i'}$ of $U$ by $N$ balls of radius $r_i \in (r_0/2, r_0)$, so that         \begin{equation}\label{eq:area-of-Bi}
            \mathcal{H}^n((\partial B_i')\cap(\Omega'\cap U)) \leq \frac{4\eps_0}{r_0} \leq \frac{\eps_1}{5N}.
        \end{equation}
        Note that the second inequality follows from the first term in the definition of $\eps_0$ in~\eqref{eq:value-of-eps0}: since $\eps_0 \leq \frac{\eps_1 r_0}{20N}$, we have $\frac{4\eps_0}{r_0} \leq \frac{4\eps_1 r_0}{20Nr_0} = \frac{\eps_1}{5N}$.

        We are going to inductively push $\Omega'$ outside of $B_1'\cap U, B_2'\cap U ,\dots, B_N'\cap U$:
        \begin{figure}[htbp]
            \centering
            \begin{tikzpicture}[scale=0.7]
              % U region (pink/salmon, hill shape on the right)
              \fill[red!20]
                (2,0) .. controls (2,1.5) and (3,2.5) .. (4,3)
                .. controls (6,4) and (8,5) .. (10,4)
                .. controls (11,3.5) and (11,1) .. (11,0) -- cycle;

              % Omega' region (sage green, left block, semi-transparent so red shows through)
              \fill[green!40!gray, fill opacity=0.5]
                (-1,0) -- (-1,6) -- (3,6)
                .. controls (3,5) and (3.5,4) .. (4,3)
                .. controls (4.5,2) and (5,1) .. (5,0) -- cycle;

              % Hatching: Omega' outside circle (left side of blue line)
              \begin{scope}
                \clip (-1,0) -- (-1,6) -- (3,6)
                  .. controls (3,5) and (3.5,4) .. (4,3)
                  .. controls (4.5,2) and (5,1) .. (5,0) -- cycle;
                \begin{scope}
                  \pgfseteorule
                  \clip (-2,-1) rectangle (12,7) (4,3) circle (1.2);
                  \fill[pattern=north east lines, pattern color=cyan!60!blue]
                    (-2,-1) rectangle (12,7);
                \end{scope}
              \end{scope}

              % Hatching: inside circle, inside Omega', outside U
              \begin{scope}
                \clip (4,3) circle (1.2);
                \clip (-1,0) -- (-1,6) -- (3,6)
                  .. controls (3,5) and (3.5,4) .. (4,3)
                  .. controls (4.5,2) and (5,1) .. (5,0) -- cycle;
                \begin{scope}
                  \pgfseteorule
                  \clip (-2,-1) rectangle (12,7)
                    (2,0) .. controls (2,1.5) and (3,2.5) .. (4,3)
                    .. controls (6,4) and (8,5) .. (10,4)
                    .. controls (11,3.5) and (11,1) .. (11,0) -- cycle;
                  \fill[pattern=north east lines, pattern color=cyan!60!blue]
                    (-2,-1) rectangle (12,7);
                \end{scope}
              \end{scope}

              % Boundary of U (red)
              \draw[very thick, red!70!black]
                (2,0) .. controls (2,1.5) and (3,2.5) .. (4,3)
                .. controls (6,4) and (8,5) .. (10,4)
                .. controls (11,3.5) and (11,1) .. (11,0);

              % Boundary of Omega' - lower part (green)
              \draw[very thick, green!50!black]
                (4,3) .. controls (4.5,2) and (5,1) .. (5,0);

              % Boundary of Omega' - upper part (blue, to intersection point)
              \draw[ultra thick, blue!70]
                (3,6) .. controls (3,5) and (3.5,4) .. (4,3);

              % Omega' lower part outside circle (blue overdraw)
              \begin{scope}
                \pgfseteorule
                \clip (-2,-1) rectangle (12,7) (4,3) circle (1.2);
                \draw[ultra thick, blue!70]
                  (4,3) .. controls (4.5,2) and (5,1) .. (5,0);
              \end{scope}

              % U boundary left part inside circle (blue)
              \begin{scope}
                \clip (4,3) circle (1.2);
                \draw[ultra thick, blue!70]
                  (2,0) .. controls (2,1.5) and (3,2.5) .. (4,3);
              \end{scope}



              % Circle B_i' (brown)
              \draw[thick, brown!70!black] (4,3) circle (1.2);

              % Circle arc inside intersection of Omega' and U (blue)
              \begin{scope}
                \clip (2,0) .. controls (2,1.5) and (3,2.5) .. (4,3)
                  .. controls (6,4) and (8,5) .. (10,4)
                  .. controls (11,3.5) and (11,1) .. (11,0) -- cycle;
                \clip (-1,0) -- (-1,6) -- (3,6)
                  .. controls (3,5) and (3.5,4) .. (4,3)
                  .. controls (4.5,2) and (5,1) .. (5,0) -- cycle;
                \draw[ultra thick, blue!70] (4,3) circle (1.2);
              \end{scope}
              \node[brown!70!black, above right] at (4.8,3.8) {$B_i'$};

              % Bottom line (partial M)
              \draw[very thick, brown!70!black] (-1,0) -- (11,0);

              % Labels
              \node[green!50!black, font=\large] at (1.5,3.5) {$\Omega'$};
              \node[red!50!black, font=\large] at (7.5,3) {$\mathcal{U}$};
              \node[blue!70, above] at (1.5,6) {$\Omega'\setminus(B_i'\cap \mathcal{U})$};
              \node[brown!70!black, right] at (11,-0.4) {\color{brown!70!black}$\partial M$};
            \end{tikzpicture}
            \caption{Pushing $\Omega'$ outside of $B_1'\cap U$. The region $\Omega'$ (green) overlaps with $\mathcal{U}$ (pink) near the ball $B_i'$. The set $\Omega' \setminus (B_i' \cap \mathcal{U})$ is obtained by removing the portion of $\Omega'$ inside $B_i' \cap \mathcal{U}$.}
            \label{fig:2675}
        \end{figure}
        
        We can start with pushing $\Omega'$ outside of $B_1'\cap U$ as in Figure~\ref{fig:2675}. By Lemma~\ref{lem:morse-separate}, for any $\eps > 0$ sufficiently small (smaller than $\frac{\eps_1}{20N}$ in fact), we can construct a nested family $\br{{\Xi_t^1}}$ that starts on $\Omega'$ and ends on a smoothing of $\Omega'\backslash(B_1'\cap U)$. We can choose the smoothing so that the set does not intersect $B_1'\cap U$. And in this step, 
        \begin{equation*}
            \begin{split}
                \mathbf{M}(\partial_I\Xi_t^1) \leq& \mathbf{M}(\partial_I\Omega') + \mathbf{M}(\partial_I(B_1'\cap U\cap \Omega'))\\
                &+2^nC(n)\mathcal{H}^{n+1}(B_1'\cap U\cap \Omega') + \eps\\
                \leq& \mathbf{M}(\partial_I\Omega') + \mathcal{H}^n(B_1'\cap \partial_I U) + \mathcal{H}^n((\partial B_1')\cap (\Omega'\cap U))\\
                &+2^nC(n)\eps_0^{\frac{n}{n + 1}} + \eps \quad\text{By \ref{eq:value-of-eps0}}\\
                \leq& \mathbf{M}(\partial_I\Omega') + \mathcal{H}^n(B_1'\cap \partial_I U) + \frac{\eps_1}{5N} + \frac{\eps_1}{4N} + \eps\quad\text{by~\ref{eq:area-of-Bi}}\\
                \leq& \mathbf{M}(\partial_I\Omega') + \mathcal{H}^n(B_1'\cap \partial_I U) + \frac{\eps_1}{2N}
            \end{split}
        \end{equation*}
        We can iterate this procedure for each ball $B_i'$. During the $i$-th step, $1 \leq i \leq N$, we have that 
        \begin{equation*}
            \mathbf{M}(\partial_I\Xi_t^i) \leq \mathbf{M}(\partial_I\Omega') + \mathcal{H}^n( \partial_I U) + \frac{i\eps_1}{2N}
        \end{equation*}
        That means we can conclude (3)(A) after gluing these together via Proposition~\ref{prop:gluing-two-nested}.
    \end{proof}
\end{lemma}

\section{Proof of the Main Result by Contradiction}\label{sec:proof-by-contradiction}

We are now ready to prove Proposition~\ref{prop:non-escape} by contradiction. Suppose there exists
a good sweepout $\br{\Gamma_t}_{t \in [0, 1]}$ such that whenever
\begin{equation*}
    \mathcal{H}^n(\Gamma_t) \geq W_g(U)
\end{equation*}
we have $\mathcal{H}^n(\Gamma_t\cap U) < \eps(U)$. Our goal is to construct a good sweepout entirely
inside $U$ whose maximal slice area is \emph{strictly less} than $W_g(U)$, contradicting the
definition of $W_g(U)$ as the infimum of maximal areas over all good sweepouts of $U$.

\subsection{Continuity Regularization}

Let $f(t) = \mathcal{H}^n(\Gamma_t\cap U)$. Note that $f(t)$ may not be continuous. However, it is easy
to see that one can perturb the family $\br{\Gamma_t}$ so that it is $\delta$-continuous in the
following sense.
\begin{definition}
    Function $f(t)$ is \textbf{$\delta$-continuous} if the oscillation
    \begin{equation*}
        \omega_f(t) := \lim_{a \to 0}[\sup_{s \in [t - a, t + a]}f(s) - \inf_{s \in [t - a, t + a]}f(s)]
    \end{equation*}
    satisfies $\omega_f(t) < \delta$ for every $t$.
\end{definition}
\begin{lemma}[$\delta$-continuous perturbation]\label{lem:delta-cts-perturbation}
    Let $U$ be a bounded relative open subset to $M$ with smooth relative boundary and $\br{\Gamma_t}$ be a good sweepout of $U$. For every $\delta > 0$ there exists a good sweepout $\br{\Gamma'_t}$ of $U$, such that 
    \begin{itemize}
        \item $f(t) = \mathcal{H}^n(\Gamma_t'\cap U)$ is $\delta$-continuous;
        \item $\sup_{t}\mathcal{H}^n(\Gamma_t') \leq \sup_{t}\mathcal{H}^n(\Gamma_t) + \delta$;
        \item $\sup_{t}\mathcal{H}^n(\Gamma_t'\cap U) \leq \sup_{t}\mathcal{H}^n(\Gamma_t\cap U) + \delta$.
    \end{itemize}
    \begin{proof}
        Fix $\delta>0$ and set $\varepsilon:=\delta/6$.
        By Lemma~\ref{lem:preliminary-modification} with parameter $\varepsilon$, from $\{\Gamma_t\}$ we obtain a piecewise nested family $\{\widetilde{\Gamma}_t\}$ and a partition $0=t_0<\cdots<t_N=1$ such that
        \begin{equation*}
        \begin{split}
        \sup_{t}\mathcal{H}^n(\widetilde{\Gamma}_t)
        &\le \sup_{t}\mathcal{H}^n(\Gamma_t)+\varepsilon,\\
        \sup_{t}\mathcal{H}^n(\widetilde{\Gamma}_t\cap U)
        &\le \sup_{t}\mathcal{H}^n(\Gamma_t\cap U)+\varepsilon.
        \end{split}
        \end{equation*}
        Indeed, the preliminary modification acts only inside finitely many thin annuli and endpoint collars; by (iii)–(iv) of Lemma~\ref{lem:preliminary-modification} the $M$–area increase is $\le\varepsilon$, and restricting to $U$ with the relative coarea formula (Corollary~\ref{cor:coarea-filling}) yields the same $\varepsilon$–bound for the $U$–area.
    
        For each subinterval $[t_i,t_{i+1}]$ given by the preliminary modification, we may represent
        $\br{\widetilde{\Gamma}_t}_{t\in[t_i,t_{i+1}]}$ as level sets of an admissible Morse function.
        Let $\mathcal{S}$ denote the finite set consisting of the partition points $\{t_1,\dots,t_{N-1}\}$ together
        with all critical values of these Morse functions. Choose $\eta>0$ so small that the open intervals
        \begin{equation*}
            J_s:=(s-\eta,s+\eta)\subset(0,1),\qquad s\in\mathcal{S}
        \end{equation*}
        are pairwise disjoint. We keep the family unchanged outside the buffers and set
        \begin{equation*}
            \Gamma_t':=\widetilde{\Gamma}_t\qquad\text{for all }t\notin \bigcup_{s\in\mathcal S}J_s.
        \end{equation*}
        All subsequent modifications will be supported inside the union of the buffers $\bigcup_{s\in\mathcal S}J_s$.
        
        Fix $s\in\mathcal S$ and set
        \begin{equation*}
            \Omega^-:=\widetilde{\Omega}_{s-\eta}\subset \Omega^+:=\widetilde{\Omega}_{s+\eta},
            \qquad
            D_s:=\Omega^+\setminus\Omega^-.
        \end{equation*}
        Shrinking $\eta>0$ if necessary and using the relative coarea formula (Corollary~\ref{cor:coarea-filling}) and the small-volume choice from Lemma~\ref{lem:preliminary-modification}, we may assume
        \begin{equation*}
          \mathcal{H}^n(\partial D_s\cap U)+
          (\mathcal{H}^{n+1}(D_s))^{\frac{n}{n+1}}
          <\varepsilon_s,
        \end{equation*}
        with $\sum_{s\in\mathcal S}\varepsilon_s\le \varepsilon$. Cover $\overline{D_s}$ by the bilipschitz balls used in the preliminary modification (interior balls mapped to balls in $\mathbb R^{n+1}$; boundary balls, via a fixed collar, to half-balls in $\mathbb R^{n+1}_+$, with vector fields tangent to $\partial M$).
        
        Applying Lemma~\ref{lem:existence-nested-between-sets} and gluing, we obtain a nested family $\{\Gamma_t'\}_{t\in J_s}$ connecting $\partial\Omega^-$ to $\partial\Omega^+$ such that, for all $t\in J_s$,
        \begin{equation*}
            \begin{split}
                \mathcal{H}^n(\Gamma_t')
                \le& \mathbf{M}(\partial_I\Omega^-)
                   + \mathbf{M}(\partial_I D_s)\\
                   &+ C(n)(1+L_{\mathrm{ball}})^n\big(\mathcal{H}^{n+1}(D_s))^{\frac{n}{n+1}}
                   + \varepsilon_s,\\
                \mathcal{H}^n(\Gamma_t'\cap U)
                \le& \mathcal{H}^n(\partial\Omega^-\cap U)
                   + \mathcal{H}^n(\partial D_s\cap U)\\
                   &+ C(n)(1+L_{\mathrm{ball}})^n\big(\mathcal{H}^{n+1}(D_s))^{\frac{n}{n+1}}
                   + \varepsilon_s,
            \end{split}
        \end{equation*}
        where $L_{\mathrm{ball}}:=(1+\frac{\varepsilon}{100W})^{1/n}-1$ so that $(1+L_{\mathrm{ball}})^n\le 1+\frac{\varepsilon}{100W}$.    
        
        By shrinking $\eta>0$ if necessary, choose numbers $\{\varepsilon_s\}_{s\in\mathcal S}$ so that the estimates in the previous paragraph hold and
        \begin{equation*}
          \sum_{s\in\mathcal S}\varepsilon_s \le \varepsilon.
        \end{equation*}
        Define the final family by
        \begin{equation*}
          \Gamma_t' :=
          \begin{cases}
            \widetilde{\Gamma}_t, & t\notin \displaystyle\bigcup_{s\in\mathcal S} J_s,\\[2pt]
            \text{the bridged slice constructed on }J_s, & t\in J_s \text{ for some } s\in\mathcal S.
          \end{cases}
        \end{equation*}
        Using the sup–bounds from the preliminary modification and the bridge estimates on each $J_s$, we obtain
        \begin{equation*}
            \begin{split}
                \sup_t \mathcal{H}^n(\Gamma_t')
                &\le \sup_t \mathcal{H}^n(\Gamma_t) + \varepsilon + \sum_{s\in\mathcal S}\varepsilon_s \le \sup_t \mathcal{H}^n(\Gamma_t)+2\varepsilon,\\
                \sup_t \mathcal{H}^n(\Gamma_t'\cap U)
                &\le \sup_t \mathcal{H}^n(\Gamma_t\cap U)+\varepsilon+\sum_{s\in\mathcal{S}}\varepsilon_s \le \sup_t \mathcal{H}^n(\Gamma_t\cap U)+2\varepsilon.
            \end{split}
        \end{equation*}
        (Any endpoint adjustment inside a fixed collar of $\partial U$ preserves the “good” endpoints and adds at most $\varepsilon$ to the above sup–bounds.)
    \end{proof}
\end{lemma}
% --- Start of the "no-escape" block with Lemma 7.4 (relative version) ---

Let $g:[0,1]\to[0,\mathcal{H}^{n+1}(U)]$ be defined by
\begin{equation*}
  g(t):=\mathcal{H}^{n+1}(U\cap\Omega_t),
\end{equation*}
which is continuous. % relative coarea; cf. CL p.25
By Lemma~\ref{lem:nested-sweepout-small-intersection} (2), if $I'$ is a connected component of
$g^{-1}([\eps_0,\ \mathcal{H}^{n+1}(U)-\eps_0])$, then for every $t\in I'$ we have
\begin{equation*}
  f(t):=\mathcal{H}^{n}(\Gamma_t\cap U) \ge 2\eps.
\end{equation*}

Applying Lemma~\ref{lem:delta-cts-perturbation} with $\delta<\eps/10$ we obtain a good sweepout
$\br{\Gamma_t'}$ such that $f'(t):=\mathcal{H}^{n}(\Gamma_t'\cap U)$ is $\delta$-continuous. 
By continuity of $g$ and the sweepout property, there exists a closed interval $I=[t_0,t_1]$ with $I'\subset I\subset[0,1]$ such that
\begin{equation}\label{eq:end-of-gt-i}
  g(t_0)\le \varepsilon_0,\qquad g(t_1)\ge \mathcal{H}^{n+1}(U)-\varepsilon_0,
\end{equation}
and (possibly shrinking $I$ slightly) we have for all $t\in I$
\begin{equation*}
  f'(t) \ge \frac{3}{2}\varepsilon.
\end{equation*}
Moreover, by the choice of endpoints and the $\delta$–continuity (with $\delta<\varepsilon/10$) together with the contrapositive of Lemma~\ref{lem:nested-sweepout-small-intersection} (2), we may assume
\begin{equation*}
  \varepsilon \le f'(t_i) \le 2\varepsilon, \qquad i=0,1.
\end{equation*}
For the upper bound, note that the contrapositive of Lemma~\ref{lem:nested-sweepout-small-intersection}\,(2) says:
if $f'(t_i)>2\varepsilon$ then 
\begin{equation*}
    g(t_i)\in\bigl(\varepsilon_0,\ \mathcal{H}^{n+1}(U)-\varepsilon_0),
\end{equation*}
which contradicts our choice of the endpoints (see~\eqref{eq:end-of-gt-i} for the bounds on $g(t_i)$).

For the lower bound, since $f'$ is $\delta$-continuous with $\delta<\varepsilon/10$ and
on the interior interval we have $f'\ge \tfrac{3}{2}\varepsilon$, choosing $t_i$ sufficiently close to the boundary of $I'$ yields
\begin{equation*}
    f'(t_i) \ge \tfrac{3}{2}\varepsilon-\delta > \varepsilon.
\end{equation*}

\begin{remark}[Why enlarge $I'$ to a closed interval $I$]
Enlarging a connected component $I'\subset B$ to $I=[t_0,t_1]$ with 
$g(t_0)\le\varepsilon_0$ and $g(t_1)\ge \mathcal{H}^{n+1}(U)-\varepsilon_0$ is essential:
(i) it yields a thin error set $P$ (hence the inward perturbation $\overline U$ with controlled relative boundary);
(ii) by $\delta$–continuity it provides endpoint anchors $\varepsilon\le f(t_i)\le 2\varepsilon$ needed for the cut–cap–insert step;
(iii) on the compact interval $I$ we obtain a strict gap $\sup_{t\in I}\mathcal{H}^n(\Gamma_t)<W_g(U)$, 
which survives the nested replacement and underwrites the final contradiction.
\end{remark}
\begin{figure}[htbp]
  \centering
  \begin{tikzpicture}[y=1cm]
    % --- Background bands (below axis) ---
    % B bands (green, matching Omega' color scheme)
    \fill[green!15!gray!20] (0.3,-0.4) rectangle (1.2,0.4);
    \fill[green!15!gray!20] (2.0,-0.4) rectangle (3.2,0.4);
    \fill[green!15!gray!20] (5.5,-0.4) rectangle (8.5,0.4);

    % A band (red, matching U color scheme)
    \fill[red!20] (3.8,-0.4) rectangle (4.8,0.4);

    % --- Axis ---
    \draw[-{Stealth}] (-0.3,0) -- (10.5,0) node[below=3pt] {$t$};
    \draw (0,0.1) -- (0,-0.1) node[below=4pt] {$0$};
    \draw (10,0.1) -- (10,-0.1) node[below=4pt] {$1$};
    \draw (5.0,0.1) -- (5.0,-0.1) node[below=4pt] {$t_0$};
    \draw (9.0,0.1) -- (9.0,-0.1) node[below=4pt] {$t_1$};

    % --- I = [t0,t1] ---
    \draw[line width=2.5pt, cyan!40!blue, opacity=0.5] (5.0,0) -- (9.0,0);
    % Dashed lines from t0, t1 down to brace
    \draw[dashed, gray] (5.0,-0.1) -- (5.0,-1.0);
    \draw[dashed, gray] (9.0,-0.1) -- (9.0,-1.0);
    \draw[decorate, decoration={brace, amplitude=5pt, mirror}]
      (5.0,-1.1) -- (9.0,-1.1) node[midway, below=6pt] {$I=[t_0,t_1]$};

    % --- I' (connected component of B inside I) ---
    \draw[line width=2pt, cyan!80!blue] (5.5,0) -- (8.5,0);
    \node[cyan!80!blue] at (7.0,0.55) {$I'$};

    % --- Conditions with arrows ---
    \node at (3.0,1.6) {$g(t_0)\le \varepsilon_0$};
    \draw[-{Stealth}, thick] (4.0,1.4) -- (5.0,0.15);

    \node at (9.0,1.6) {$g(t_1)\ge \mathrm{Vol}(U)-\varepsilon_0$};
    \draw[-{Stealth}, thick] (9.0,1.3) -- (9.0,0.15);

    % --- Legend ---
    \fill[green!15!gray!20] (0.3,-2.5) rectangle (0.7,-2.2);
    \node[anchor=west, green!50!black, font=\small] at (0.8,-2.35)
      {$B=\{t: g(t)\in(\varepsilon_0,\,\mathrm{Vol}(U)-\varepsilon_0)\}$};
    \fill[red!20] (6.5,-2.5) rectangle (6.9,-2.2);
    \node[anchor=west, red!70!black, font=\small] at (7.0,-2.35)
      {$A=\{\mathcal{H}^n(\Gamma_t)\ge W_g(U)\}$};
  \end{tikzpicture}
  \caption{Time–axis schematic. The blue bands show $B=\{t:\ g(t)\in(\varepsilon_0,\ \mathrm{Vol}(U)-\varepsilon_0)\}$; the red band shows $A=\{t:\ \mathcal H^n(\Gamma_t\cap M)\ge W_g(U)\}$ with $A\cap B=\varnothing$. A connected component $I'\subset B$ is enlarged to $I=[t_0,t_1]$ with $g(t_0)\le\varepsilon_0$, $g(t_1)\ge \mathrm{Vol}(U)-\varepsilon_0$.}
  \label{fig:time-axis-schematic}
\end{figure}

Relabelling $f'$ back to $f$ and $\Gamma_t'$ back to $\Gamma_t$, define
\begin{equation*}
    \begin{split}
        A&:=\br{t\in[0,1]:\ \mathcal{H}^n(\Gamma_t)\ge W_g(U)},\\
        B&:=\br{t\in[0,1]:\ g(t)\in(\varepsilon_0,\ \mathcal{H}^{n+1}(U)-\varepsilon_0)}.
    \end{split}
\end{equation*}
By Lemma~\ref{lem:nested-sweepout-small-intersection} (2), $t\in B\Rightarrow f(t)\ge 2\varepsilon$, hence $A\cap B=\emptyset$. Since $I\subset B$, we have
\begin{equation*}
\mathcal{H}^n(\Gamma_t)\ <\ W_g(U)\qquad\text{for all }t\in I.
\end{equation*}
By piecewise continuity of $t\mapsto \mathcal{H}^n(\Gamma_t)$ on the compact interval $I$, there exists
\begin{equation*}
m:=\sup_{t\in I}\mathcal{H}^n(\Gamma_t) < W_g(U),\qquad
\delta_1:=W_g(U)-m > 0.
\end{equation*}

Apply the nested–replacement proposition from Section~6 (Proposition~\ref{prop:nested-sweepout}) to $\br{\Gamma_t}_{t\in I}$. 
We obtain a nested family $\{\overline{\Gamma}_t\}_{t\in[0,1]}$ with associated open sets $\{\overline{\Omega}_t\}$ such that
\begin{equation*}
  \sup_{t\in[0,1]}\mathcal{H}^n(\overline{\Gamma}_t)\ \le\ m-\tfrac{\delta_1}{2}\ <\ W_g(U)-\tfrac{\delta_1}{2},
\end{equation*}
and
\begin{equation*}
  \overline{\Omega}_0\subset \Omega_{t_0},\qquad \Omega_{t_1}\subset \overline{\Omega}_1.
\end{equation*}

Define the thin error set
\begin{equation*}
  P \ :=\ (\Omega_{t_0}\cap U)\ \cup\ \bigl(U\setminus \operatorname{cl}(\Omega_{t_1})).
\end{equation*}
By the endpoint volume bounds $g(t_0)\le \varepsilon_0$ and $g(t_1)\ge \mathcal{H}^{n+1}(U)-\varepsilon_0$, both
$\Omega_{t_0}\cap U$ and $U\setminus \operatorname{cl}(\Omega_{t_1})$ have measure $\le \varepsilon_0$, hence
$\mathcal{H}^{n+1}(P)\le 2\varepsilon_0$. Remove $P$ from $U$ and smooth inward inside a fixed collar of $\partial U$
to obtain a slightly shrunken open set $\overline U\subset U$ (an relative inward $\delta$-perturbation of $U\setminus \operatorname{cl}(P)$).
Then, as in \cite[p.~26]{CL20}, we may assume:
\begin{enumerate}
  \item[(i)] \emph{$\{\overline{\Gamma}_t\}$ restricts to a nested sweepout of $\overline U$.}
  Since $\overline\Omega_0\subset \Omega_{t_0}$ and $\Omega_{t_1}\subset \overline\Omega_1$, removing the very thin
  set $P$ does not break nestedness; reparametrising the family and performing corner-smoothing inside the collar of
  $\partial U$ yields a nested sweepout of $\overline U$.
  \item[(ii)] \emph{Relative boundary control:}
  \begin{equation*}
    \mathbf{M}(\partial_I\overline U)
    \ \le\ \mathbf{M}(\partial_IU)\ +\ C\,\varepsilon_0^{\frac{n}{n+1}}
    \ \le\ 2\,\mathbf{M}(\partial_IU)
  \end{equation*}
  for $\varepsilon_0=\varepsilon_0(U)$ chosen sufficiently small. Here the extra term comes from the new boundary
  created by excising $P$ and the inward smoothing; it is controlled by the relative isoperimetric/coarea
  inequalities applied to $\mathcal{H}^{n+1}(P)\le 2\varepsilon_0$.
  \item[(iii)] \emph{Push-out with controlled area:}
  there exists a homotopy supported in the collar of $\partial U$, generated by vector fields tangent to $\partial M$,
  that pushes $\partial_I\overline U$ outside $U$ through slices of area at most $3\mathbf{M}(\partial_IU)$, after absorbing the $O(\varepsilon_0^{\frac{n}{n+1}})$ error into the constant (again by choosing
  $\varepsilon_0(U)$ small). This is the relative version of the short “push across the collar” used in \cite{CL20}.
\end{enumerate}
\subsection{Controlling Slices Outside the Good Set}

Before proceeding, we isolate a slice of the nested family whose portion \emph{outside} the shrunken
set $\overline U$ is uniformly controlled. This is the key bridge between the global width bound and
a sweepout confined to $\overline U$: once such a slice is found, we can (i) cut the family at that
time, (ii) cap the outside part by using only boundary data of $U$, and (iii) build a sweepout of
$\overline U$ whose maximal area is strictly below $W_g(U)$. The next lemma provides precisely this
slice-level estimate and makes the subsequent "cut–and–insert" step possible.
\begin{lemma}[relative version of Lemma 7.4 in \cite{CL20}]\label{lem:barU-slice-outside}
There exists $t_0\in[0,1]$ such that
\begin{equation*}
  \mathcal{H}^n(\overline{\Gamma}_{t_0}\setminus \overline U)\leq2\mathbf{M}(\partial_I U).
\end{equation*}
\end{lemma}
\begin{proof}
    Let $L:=\sup_{t\in[0,1]}\mathcal{H}^n\bigl(\overline{\Gamma}_t\cap \overline U)$. By the Morse–foliation construction on $\overline U$ and the relative boundary control in (ii), we can sweep out $\overline U$
    by slices of area at most
\begin{equation*}
  L+\mathbf{M}(\partial_I\overline U)+\delta \le L+2\,\mathbf{M}(\partial_IU)+\delta.
\end{equation*}
Using (iii), we extend this to a good sweepout of $U$ (the homotopy is supported in a fixed collar of $\partial U$ and generated by vector fields tangent to $\partial M$) whose maximal slice area is at most $L+2\,\mathbf{M}(\partial_IU)+\delta$. Hence, by the definition of $W_g(U)$,
\begin{equation*}
  W_g(U)\le L+2\mathbf{M}(\partial_IU)+\delta.
\end{equation*}
If, for every $t\in[0,1]$, we had $\mathcal{H}^n(\overline{\Gamma}_t\setminus \overline U)\ >\ 2\,\mathbf{M}(\partial_IU),$
then
\begin{equation*}
    \begin{split}
        \mathcal{H}^n(\overline{\Gamma}_t) &= \mathcal{H}^n(\overline{\Gamma}_t\cap \overline U) + \mathcal{H}^n( (\overline{\Gamma}_t\cap M)\setminus \overline U)\\
        &> L + 2\mathbf{M}(\partial_IU)\\
        &\geq W_g(U)-\delta
    \end{split}
\end{equation*}
contradicting $\sup_t\mathcal{H}^n(\overline{\Gamma}_t)\le W_g(U)-\delta_1/2$ for $\delta>0$ chosen smaller than $\delta_1/2$.
Therefore such $t_0$ exists.
\end{proof}

Let 
\begin{equation*}
    \begin{split}
        &\mathcal{U}_0:=\br{\Omega: \overline{\Omega}_0\subset \Omega\subset \overline{\Omega}_{t_0}\setminus \overline{U}};\\
        &\mathcal{U}_1:=\br{\Omega: \overline{\Omega}_{t_0}\cup \overline{U}\subset \Omega\subset \overline{\Omega}_1}
    \end{split}
\end{equation*}
and define
\begin{equation*}
  A_i:=\inf\br{\mathbf{M}(\partial_I\Omega) : \Omega\in\mathcal{U}_i},\qquad i=0,1.
\end{equation*}
By Lemma~\ref{lem:barU-slice-outside} and property (ii) for $\overline U$, we have 
\begin{equation*}
  A_i\ \le\ 3\,\mathbf{M}(\partial_IU).
\end{equation*}
Choose $\Sigma_i=\partial\Xi_i$ with $\Xi_i\in\mathcal U_i$ such that 
\begin{equation*}
  \mathcal{H}^n(\Sigma_i)\ \le\ A_i+\delta/4.
\end{equation*}
Applying the relative nested–insertion lemmas (Lemma~\ref{lem:extension}(I),(II) in our Section~5) \emph{inside} $\overline U$ and concatenating,
we obtain a nested sweepout of $\overline U$ from $\Sigma_0$ to $\Sigma_1$ whose maximal slice area (measured in $M$) satisfies
\begin{equation*}
  \sup_{t}\ \mathcal{H}^n(\Gamma^{\overline U}_t)\ \le\ W_g(U)-\delta/4.
\end{equation*}
Using Lemma~\ref{lem:nested-sweepout-small-intersection}\,(3) (relative push across a $\partial U$–collar),
we upgrade this to a good sweepout of $U$ with
\begin{equation*}
  \sup_{t}\ \mathcal{H}^n(\Gamma^{U}_t)\ \le\ W_g(U)-\delta/4,
\end{equation*}
contradicting the definition of $W_g(U)$. This completes the proof of Proposition~\ref{prop:non-escape}.
